%% kapa: exam=2 round=22 year=2011 period=3 subject=compensation_law subject_ko=감정평가 및 보상법규 seq=2 number=2 points=30
%% source: Q-Net 감정평가사 기출문제 (cst00309) artlSeq=5040389 file=extracted/1008/5040389/2011년 제22회 감정평가사 2차시험 기출문제/3교시.hwp
%% conversion: kordoc 4.24.51 -> tools/gen2.py (auto); review=proofread-sonnet; flags=none
\begin{longquestion}{【문제 2】}
다음 각각의 사례에 대하여 답하시오. (30점) \par
\medskip\noindent (1) 국토해양부장관은 감정평가업자 甲에 대하여 법령상 의무 위반을 이유로 6월의 업무정지처분을 하였다. 甲은 업무정지처분취소소송을 제기하였으나 기각되었고 동 기각판결은 확정되었다. 이에 甲은 위 처분의 위법을 계속 주장하면서 이로 인한 재산상 손해에 대해 국가배상청구소송을 제기하였다. 이 경우 업무정지처분취소소송의 위법성 판단과 국가배상청구소송의 위법성 판단의 관계를 설명하시오. (20점) \par
\medskip\noindent (2) 감정평가업자 乙은 국토해양부장관에게 감정평가사 갱신등록을 신청하였으나 거부당하였다. 그런데 乙은 갱신등록거부처분에 앞서 거부사유와 법적 근거, 의견제출의 가능성 등을 통지받지 못하였다. 위 갱신등록거부처분의 위법성 여부를 검토하시오. (10점) \par
\end{longquestion}
